% ============================================================
% 04-programsko-rjesenje.tex - PRVA VERZIJA
% ============================================================

Programsko rje\v{s}enje razvijeno je u programskom jeziku Python, a
obuhva\'{c}a samostalnu implementaciju svih jedanaest SMOTE algoritama,
evaluacijski okvir za eksperimentalnu analizu s osam klasifikatora i
vi\v{s}e metrika evaluacije, modul za statisti\v{c}ku obradu rezultata
te web su\v{c}elje za interaktivnu vizualizaciju. U nastavku se
detaljno opisuju arhitektura i na\v{c}in rada.

\subsection{Na\v{c}in rada programskog rje\v{s}enja}

Arhitektura rje\v{s}enja sastoji se od nekoliko me\dj{}usobno povezanih
modula. Jezgru \v{c}ini paket \texttt{smote\_variants} s apstraktnom
baznom klasom \texttt{BaseSMOTE} koja definira su\v{c}elje
\texttt{fit\_resample(X, y)} \v{s}to ga svaka izvedenica mora
implementirati. Ovo omogu\'{c}uje uniformno kori\v{s}tenje svih
algoritama u evaluacijskom okviru. Implementirane su sljede\'{c}e
varijante: SMOTE, Borderline-SMOTE (obje ina\v{c}ice --- BS1 i BS2),
ADASYN, Safe-Level-SMOTE, K-Means SMOTE, SVM-SMOTE, SMOTE-ENN,
SMOTE-Tomek, Geometric SMOTE (G-SMOTE), Random SMOTE i Polynom-Fit SMOTE.
Svi algoritmi implementirani su samostalno (from-scratch), oslanjaju\'{c}i
se na biblioteke \texttt{numpy} za numeri\v{c}ke operacije i
\texttt{scikit-learn} \cite{pedregosa2011} isklju\v{c}ivo za KNN
pretra\v{z}ivanje (\texttt{NearestNeighbors}), klasteriranje
(\texttt{KMeans}) i SVM (\texttt{SVC}) kod odgovaraju\'{c}ih izvedenica.

Za svaki SMOTE algoritam predvi\dj{}eno je iterativno ispitivanje vi\v{s}e
vrijednosti klju\v{c}nih parametara --- primjerice, broj susjeda $k$
testirat \'{c}e se u rasponu $\{3, 5, 7, 10\}$ --- kako bi se za svaki
skup podataka i klasifikator identificirala konfiguracija koja daje
najbolje rezultate. Ovaj pristup omogu\'{c}uje pravednu usporedbu
algoritama jer svaki od njih dolazi do izra\v{z}aja u svojoj optimalnoj
postavci, \v{c}ime se izbjegava pristranost koja bi nastala fiksiranjem
istih vrijednosti parametara za sve varijante.

Evaluacijski modul implementira stratificiranu $k$-struku unakrsnu
validaciju ($k = 5$) za svaku kombinaciju algoritma preuzorkovanja i
klasifikatora. Kori\v{s}teni klasifikatori uklju\v{c}uju stablo odluke
(DT), slu\v{c}ajnu \v{s}umu (RF), XGBoost \v{c}ija je prednost ugra\dj{}ena
podr\v{s}ka za neuravnote\v{z}ene klase kroz parametar
\texttt{scale\_pos\_weight}, logisti\v{c}ku regresiju (LR), SVM s RBF
jezgrom, $k$-NN, Gaussov naivni Bayes (GNB) te vi\v{s}eslojni perceptron
(MLP) s jednim skrivenim slojem kao jednostavna neuronska referentna
to\v{c}ka. Svi klasifikatori preuzeti su iz biblioteke scikit-learn, uz
iznimku XGBoosta koji dolazi iz istoimene biblioteke.

Za evaluaciju performansi prikupljat \'{c}e se \v{s}iri skup metrika
nego \v{s}to je to uobi\v{c}ajeno u sli\v{c}nim radovima, s ciljem
\v{s}to potpunijeg i znanstveno utemeljenog uvida u pona\v{s}anje
svakog algoritma. Iako su neke od ovih mjera ra\v{c}unski zahtjevnije
od jednostavne to\v{c}nosti klasifikacije --- primjerice, AUC-PR zahtijeva
izra\v{c}un krivulje u vi\v{s}e to\v{c}aka, a MCC uzima u obzir sva
\v{c}etiri elementa matrice zbunjenosti --- upravo ta slo\v{z}enost
omogu\'{c}uje pouzdaniju procjenu na neuravnote\v{z}enim podacima.
Prikupljat \'{c}e se sljede\'{c}e metrike:
F-1 mjera kao harmonijska sredina preciznosti i odziva; G-Mean kao
geometrijska sredina odziva i specifi\v{c}nosti, \v{s}to je posebno
prikladno za neuravnote\v{z}ene podatke jer ka\v{z}njava modele koji
favoriziraju jednu klasu; AUC-ROC kao mjera separabilnosti klasa neovisna
o pragu odluke; AUC-PR (povr\v{s}ina ispod Precision-Recall krivulje)
koja je informativnija od AUC-ROC kod ekstremne neuravnote\v{z}enosti;
Balanced Accuracy kao aritmeti\v{c}ka sredina odziva i specifi\v{c}nosti;
Matthewsov koeficijent korelacije (MCC) koji pru\v{z}a sveobuhvatnu ocjenu
temeljem sva \v{c}etiri elementa matrice zbunjenosti i smatra se jednom od
najpouzdanijih pojedina\v{c}nih mjera za neuravnote\v{z}ene probleme; te
F2-mjera koja daje ve\'{c}u te\v{z}inu odzivu, \v{s}to je relevantno u
domenama gdje su la\v{z}no negativni primjeri skuplji od la\v{z}no
pozitivnih.

Modul za statisti\v{c}ku analizu provodit \'{c}e tri testa \v{c}iji
\'{c}e se rezultati koristiti pri interpretaciji eksperimentalnih
nalaza u petom poglavlju. Friedmanov test \cite{demsar2006} slu\v{z}i za
utvr\dj{}ivanje postoje li statisti\v{c}ki zna\v{c}ajne razlike me\dj{}u
algoritmima na razini svih skupova podataka --- odnosno, nije li uo\v{c}ena
razlika u performansama tek posljedica slu\v{c}ajnosti. U slu\v{c}aju
odbacivanja nul-hipoteze Friedmanovog testa, Nemenyi post-hoc test
identificira koji konkretni parovi algoritama se statisti\v{c}ki
zna\v{c}ajno razlikuju, a rezultati se prikazuju kriti\v{c}nim
dijagramima. Wilcoxon signed-rank test koristit \'{c}e se za direktnu
usporedbu svake izvedenice s osnovnim SMOTE algoritmom kao referentnom
to\v{c}kom, \v{c}ime se dobiva odgovor na pitanje donosi li odre\dj{}ena
modifikacija SMOTE-a zna\v{c}ajno pobolj\v{s}anje ili ne. Svi testovi
provodit \'{c}e se kori\v{s}tenjem biblioteke \texttt{scipy.stats}.

Vizualizacijski modul slu\v{z}i za grafi\v{c}ki prikaz rezultata kroz
CD dijagrame, box-plotove, heatmap-e i scatter prikaze sinteti\v{c}kih
primjera.

\subsection{Prikaz i na\v{c}in uporabe programskog rje\v{s}enja}

Eksperimenti se pokre\'{c}u iz komandne linije naredbom
\texttt{python run\_analysis.py}, pri \v{c}emu je mogu\'{c}e konfigurirati
parametre SMOTE algoritama, odabir klasifikatora i skupova podataka.
Rezultati eksperimenata pohranjuju se u CSV formatu i prikazuju kroz
tablice i grafove. Web su\v{c}elje --- opisano u nastavku ovog poglavlja --- omogu\'{c}it
\'{c}e interaktivnu vizualizaciju procesa generiranja sinteti\v{c}kih
primjera. Za ilustraciju \'{c}e se koristiti screenshotovi glavnih
dijelova su\v{c}elja.

Eventualno, dati dijagram toka programskog rje\v{s}enja na visokoj razini
koji prikazuje tijek od u\v{c}itavanja podataka do prikaza rezultata.
Klju\v{c}ni dijelovi programskog k\^{o}da bit \'{c}e prilo\v{z}eni u
prilogu.

\subsection{Web su\v{c}elje za vizualizaciju}

% ============================================================
% 06-web-sucelje.tex — PRVA VERZIJA
% ============================================================

Razvijeno je i web su\v{c}elje za interaktivnu
vizualizaciju procesa generiranja sinteti\v{c}kih primjera razli\v{c}itim
SMOTE varijantama. Su\v{c}elje omogu\'{c}uje korisniku da na intuitivan
na\v{c}in istra\v{z}i razlike me\dj{}u algoritmima i utjecaj njihovih
parametara na distribuciju generiranih uzoraka.

\subsubsection{Motivacija}

Razumijevanje na\v{c}ina na koji razli\v{c}iti SMOTE algoritmi generiraju
sinteti\v{c}ke primjere \v{c}esto je ote\v{z}ano zbog vi\v{s}edimenzionalne
prirode podataka. Vizualizacija u dvije ili tri dimenzije, dobivena
redukcijom dimenzionalnosti, omogu\'{c}uje intuitivno sagledavanje
razlika izme\dj{}u linearnih pristupa (klasi\v{c}ni SMOTE) i
geometrijskih pro\v{s}irenja (G-SMOTE), \v{s}to ima i zna\v{c}ajnu
edukacijsku vrijednost.

\subsubsection{Arhitektura web su\v{c}elja}

Su\v{c}elje je realizirano kori\v{s}tenjem biblioteke Streamlit za
Python koja omogu\'{c}uje brzo prototipiranje interaktivnih web
aplikacija. Backend se izravno oslanja na implementirane SMOTE algoritme
iz paketa \texttt{smote\_variants}. Za vizualizaciju se koristi Plotly
(interaktivni scatter plotovi), dok se redukcija dimenzionalnosti provodi
algoritmima PCA i t-SNE iz biblioteke scikit-learn.

\subsubsection{Funkcionalnosti}

Opisuju se klju\v{c}ne funkcionalnosti su\v{c}elja:
u\v{c}itavanje vlastitog CSV skupa podataka ili odabir jednog od
ugra\dj{}enih skupova; redukcija dimenzionalnosti na 2D/3D pomo\'{c}u
PCA ili t-SNE; odabir jedne ili vi\v{s}e SMOTE varijanti za usporedbu;
pode\v{s}avanje parametara ($k$, postotak preuzorkovanja); usporedni
prikaz originalnih i sinteti\v{c}kih primjera; animacija procesa
generiranja za 2D podatke; prikaz granice odluke klasifikatora prije i
poslije preuzorkovanja te tablica metrika za kvantitativnu usporedbu.

\subsubsection{Prikazi web su\v{c}elja}

Prikazat \'{c}e se screenshotovi glavnih ekrana su\v{c}elja s kratkim
opisima: po\v{c}etni ekran s kontrolama, PCA prikaz originalnih podataka,
usporedni prikaz vi\v{s}e SMOTE varijanti, prikaz generiranja G-SMOTE u
odnosu na osnovni SMOTE, granica odluke prije i poslije preuzorkovanja
te tablica metrika.

